\cotefile{16}
\scanpage{3}\mspage{1}

\uline{Sur certains anneaux gradués}

\inserted{Soit $k$ anneau (\unsure{comm.|commut.}) fixé.}

1. Soi\deleted{en}t \deleted{$k$ un anneau,} $M^*$ un $k$-module gradué \inserted{par $\mathbb Z$} à graduation \inserted{finie}, $\unsure{E|\mathcal E}^*$ un\deleted{e} \deleted{algèbre graduée} \inserted{sous-anneau de l'anneau} d'endomorphismes \inserted{du module \unsure{ss-jacent} (\unsure{s/s graduation})} de $M^*$. On dénote par $\pi_i$ la projection canonique de $M^*$ sur $M^i$. \deleted{et} On suppose que l'on a
\[
 \pi_i\in E^* \qquad \text{pour tout }i
 \tag*{(\unsure{1.2|1.1})}
\]
Cela implique que $E^*$ est une \uline{sous}-\deleted{\uline{algèbre}} \inserted{anneau} graduée de l'algèbre graduée \deleted{\unsure{des}} $\operatorname{End}(M^*)$ [\deleted{\unsure{dont la}} graduation est définie, rappelons le, par
\[
 \operatorname{End}^{(\unsure{\nu|v})}(M)=\sum\operatorname{Hom}(M^i,M^{i+\nu})
\]
ou de façon équivalente
\[
 \operatorname{End}^{(\nu)}(M)=\sum_i\pi_{i+\nu}\operatorname{End}(M^*)\pi_i,
\]
la composante homogène de degré $\nu$ d'un endomorphisme $u$ de $M^*$ étant
\[
 \deleted{\unsure{u|\mathrm{gr}}}\operatorname{pr}_{\nu}(u)=\sum_i\pi_{i+\nu}u\pi_i .
\]

La graduation de $E^*$ est définie \inserted{intrinsèquement}, en termes de la \deleted{\unsure{gr}} donnée des $\pi_i$, par
\[
\left\{\begin{aligned}
 \text{\textcircled{\scriptsize a}}\quad &\boxed{E^{(\nu)}=\sum_i\pi_{i+\nu}E^*\pi_i}\\
 \text{\textcircled{\scriptsize b}}\quad &\operatorname{pr}^{E}_{\nu}u=\sum_i\pi_{i+\nu}u\pi_i
\end{aligned}\right.
\tag*{(\deleted{\unsure{1.2}}1.3)}
\]
et
\[
 u\in E^{(\nu)}\iff u\pi_i=\pi_{i+\nu}u\quad\text{pour tt } i
 \tag*{(1.4.)}
\]
\deleted{\illegible{2}}
\inserted{Noter les relations
\[
 \deleted{\text{Notons}}\quad
 \left|\begin{aligned}
 \pi_i^2&=\pi_i\\
 \pi_i\pi_j&=0\quad\text{si }i\ne j\\
 \sum\pi_i&=1
 \end{aligned}\right.
 \tag*{(1.1)}
\]}

2. Soit $E^*$ un \deleted{$k$-algèbre} \inserted{anneau} gradué, et soient $\pi_i$ ($i\in\mathbb Z$) des éléments de $E$, \inserted{nuls sauf un nb fini.} A quelle condition les formules (\unsure{1.3}) définissent-elles
% archive sheet number: 2
